% !PRINTED 383
\weberpage{383}

%% [running head, not body text -- commented out 2026-08-07] Page furniture (running
%% head, folio, signature) is excluded from this edition by standing policy.
%% \noindent\S.\,124. \hfill \textsc{Reducirte Zahlen.} \hfill 383

%<<CONT
(2) $\xi_1^2 = \tfrac{1}{4}$, $\xi_1 = \mp\tfrac{1}{2}$, was sowohl unter dem Fall a) als unter dem
Fall b) enthalten ist;

und hiermit ist also der Satz 2. bewiesen. Dem fügen wir nun noch als
dritten Satz hinzu:

\noindent 3.\quad \gesperrt{Zu einer gegebenen negativen Discriminante $D$ giebt es nur eine endliche Anzahl reducirter Irrationalzahlen.}

Aus der Darstellung der Irrationalzahlen \S.\,122, (7):
\[
\omega = \dfrac{\sqrt{D}+b}{2c}, \qquad \xi = \dfrac{b}{2c}, \qquad \eta = \dfrac{\sqrt{-D}}{2c}
\]
und aus der Ungleichung (3) $\eta^2 \geqq \tfrac{3}{4}$ folgt
\[
\dfrac{-D}{4c^2} \geqq \tfrac{3}{4},
\]
oder
\[
c \leqq \sqrt{\dfrac{-D}{3}};
\]
also giebt es, wenn $D$ gegeben ist, nur eine endliche Zahl von Werthen
der positiven ganzen Zahl $c$. Ferner folgt aus $-\tfrac{1}{2} \leqq \xi \leqq \tfrac{1}{2}$:
\[
-c \leqq b \leqq c.
\]
Also giebt es zu jedem Werth von $c$ nur eine endliche Zahl von Werthen
für $b$, w.\,z.\,b.\,w.

Vereinigt man alle unter einander äquivalenten Zahlen zu einer
Zahlclasse, so haben wir dann also nach 1. bewiesen:

\noindent 4.\quad \gesperrt{Die Anzahl der Classen von imaginären quadratischen Irrationalzahlen einer gegebenen Discriminante ist endlich.}

\begin{center}
\S.\,124.

\gesperrt{Reducirte Zahlen mit positiver Discriminante.}
\end{center}

Wir haben nun eine ähnliche Untersuchung durchzuführen für ein
positives $d$, also für \gesperrt{reelle quadratische Irrationalzahlen}. Eine solche
Zahl $\omega = x + y\sqrt{d}$ soll \gesperrt{reducirt} heissen, wenn sie folgenden
Bedingungen genügt:
%CONT>>
